% =============================================================
% 04_cross_review.tex  第4部 分野横断の整理
% =============================================================
\chapter{分野横断の整理：3 分野をつなげて見る}

\section{なぜ横断で学ぶのか}

ITパスポートで失点が多いのは、実は「分野を横断する問題」です。
「セキュリティ事故が起きたときに、\textbf{経営者が負うべき責任}は？」のように、
3 分野の知識を同時に使う問題が、近年では毎回のように出題されています。

\pitfall{「一つの分野だけを完璧にする」作戦は危険}{
ストラテジを得意にしても、テクノロジの用語に弱いと横断問題で止まります。
本章では、\textbf{章同士のつながり}を一段高い視点で整理します。
}

\section{頻出の横断テーマ 3 つ}

\subsection{1. セキュリティ × 法務 × 運用}

セキュリティ事故は、技術の問題のように見えて、実は\termtag{組織のルール}と\termtag{法律上の責任}の問題でもあります。

\comparisonbx{インシデント対応で関係する知識}{%
\comptable{視点 & 使う知識 & 章との対応}{
技術面 & 攻撃手法・対策技術 & テクノロジ系（第3部） \\
運用面 & インシデント対応手順 & マネジメント系（第2部） \\
法務面 & 個人情報保護法・通知義務 & ストラテジ系（第1部）
}
}

\subsection{2. システム開発 × 契約 × 見積}

「誰が・いくらで・いつまでに」は、開発工程（マネジメント）と契約（ストラテジ）の両方にまたがります。

\definitionbx{準委任契約と請負契約}{
\termtag{請負契約}：\textbf{完成責任}あり。「動くものを作る」ことが納品物。
\termtag{準委任契約}：\textbf{善管注意義務}のみ。「決められた作業をちゃんとやる」ことが納品物。
アジャイル開発では準委任契約が選ばれやすい。
}

\subsection{3. データ活用 × ビジネス戦略}

AI・ビッグデータを扱う問題は、\textbf{技術そのもの}よりも\textbf{どう使うか}が問われる傾向にあります。

\advice{%
「DX（デジタルトランスフォーメーション）」系の問題は、
\textbf{「技術の話ではなく、経営の話」}だと意識すると、選択肢の見え方が変わります。
}

\section{ひっかけ選択肢の見抜き方}

ITパスポートの選択肢には、よく使われる\textbf{ひっかけの型}があります。

\comparisonbx{典型的なひっかけパターン}{%
\comptable{型 & 見抜き方 & 例}{
用語すり替え & 似た用語の定義をすり替える & 認証⇄認可 \\
範囲のすり替え & 一部の事例を全体のように書く & 「すべての」「必ず」 \\
手段と目的のすり替え & 目的として手段を書く & 「セキュリティの目的は暗号化である」 \\
層の間違い & 階層・分野を一段ずらす & OSI の層番号ずれ
}
}

\pitfall{「もっともらしい日本語」に負けない}{
文としてきれいでも、\textbf{用語の定義と違うことを言っているだけ}のことがあります。
読んだときに「そうだよね」と思った選択肢ほど、もう一度定義に照らし合わせましょう。
}

\onelinesummary{%
横断問題は\textbf{「どの分野の知識を使うか」}を仕分ける力が問われる。
分野別の章で覚えた知識を、\termtag{テーマ}の単位で束ね直すのが本章の目的。
}
